% Source-derived chapter 02: Notation
% Original source: chabauty-coleman-bound-surfaces
\chapter{Notation}
\label{chabauty-coleman-bound-surfaces-chapter-02}
\source{chabauty-coleman-bound-surfaces}

\begin{remark}[Source section]
\label{chabauty-coleman-bound-surfaces-chapter-02-source}
\source{chabauty-coleman-bound-surfaces}
\source{chabauty-coleman-bound-surfaces}
This chapter reproduces the corresponding section of the retrieved
LaTeX source, including its definitions, statements, examples, and
proofs. It has not yet been translated into Lean.
\notready
\end{remark}

% The source section title was: Notation
\subsection{General notation}\label{NotationFields} $L$ is a field and $L^{alg}$ is an algebraic closure of it. The symbol $F$ will denote a field, a function on a variety, or a power series, and this will be explicitly stated.  The symbol $K$ will denote a $p$-adic field or a canonical divisor in a variety, which will be explicitly stated.  $k$ will always denote an algebraically closed field and $\kk$ will always denote a finite field. 
\source{chabauty-coleman-bound-surfaces}


We write $\N=\Z_{\ge 0}$. For $\alpha=(\alpha_1,...,\alpha_n)\in \N$ we define $\|\alpha\|=\alpha_1+...+\alpha_n$. If $\xx=(x_1,...,x_n)$ is an $n$-tuple of variables, we write $\xx^{\alpha}=x_1^{\alpha_1}\cdots x_n^{\alpha_n}$. 


The exponential function on $\R$ is $\exp$ and the logarithm on $\R_{>0}$ is $\log$.
%%
%%
\subsection{Geometry}\label{notationGeo} A variety over $L$ is a reduced $L$-scheme which is separated of finte type over $L$. We remark that we are not requiring irreducibility as part of the definition. A curve (resp. surface) over $L$ is a variety over $L$ of pure dimension $1$ (resp. 2). 
\source{chabauty-coleman-bound-surfaces}

The normalization of an integral domain $A$ is denoted by $\widetilde{A}$, and the normalization of a variety $Z$ over $k$ is denoted by $\widetilde{Z}$. 

If $X$ is a scheme and $x$ is a (schematic) point in $X$, the maximal ideal of the local ring $\Ocal_{X,x}$ is denoted by $\mfrak_{X,x}$. The completion of $\Ocal_{X,x}$ at $\mfrak_{X,x}$ is denoted by $\widehat{\Ocal}_{X,x}$. 

If $f:X\to Y$ is a morphism of schemes, $x$ is a point of $X$ and $y=f(x)$, the induced map on local rings is $f^\#_x : \Ocal_{Y,y}\to \Ocal_{X,x}$. If $X$ consists of exactly one point, we just write $f^\#$ instead of $f^\#_x$.


For an irreducible projective curve $C$ over the algebraically closed field $k$, the arithmetic genus of $C$ is denoted by $\gfrak_a(C)$ and the geometric genus is denoted by $\gfrak_g(C)$. Namely, $\gfrak_g(C)=\gfrak_a(\widetilde{C})$. 

If $L$ is a perfect field and $C$ is a geometrically irreducible projective curve over $L$, we define the arithmetic and geometric genus of $C$ by base change to $L^{alg}$.

Given a geometrically irreducible, smooth, projective surface $S$ over $L$, we have a $\Z$-valued intersection product on divisors of $S$, which we write $(D.D')$, $(D.D')_S$, or $D.D'$ for two divisors $D$ and $D'$ on $S$, and we write $D^2=(D.D)$ for a self-intersection. The first Chern number $c_1^2(S)$ of $S$ is the self-intersection number of a canonical divisor on $S$. 

%%%
%%%
\subsection{Local fields}\label{NotationLoc} In the context of local fields, the symbol $p$ denotes a prime number, $K$ is a finite extension of $\Q_p$, $R$ is the ring of integers of $K$, $\pfrak$ is the maximal ideal of $R$, $\varpi\in \pfrak$ is a uniformizer, $\kk=R/\pfrak$ is the residue field of $R$, $q=\# \kk$, $\efrak$ is the ramification exponent of $K/\Q_p$, and $\ffrak=[\kk : \F_p]$ is the residue degree. We denote by $|-|$ the absolute value on $K$ normalized by $|\varpi|=1/q$, or equivalently, $|p|=p^{-[K:\Q_p]}$.
\source{chabauty-coleman-bound-surfaces}


For $n\ge 1$ we put the following norm on $K^n$:
$$
|\vv | = \max_{1\le j\le n}|v_j|, \quad \vv =(v_1,...,v_n)\in K^n.
$$
For $r>0$ we define the balls $B_n(r)=\{\vv \in K^n : |\vv |<r\}$ and  $B_n[r]=\{\vv \in K^n : |\vv |\le r\}$.
 





%%%%%%%%%%%%%%%%%%%%%%%%%%%%%%%%%%%%%%
%%%%%%%%%%%%%%%%%%%%%%%%%%%%%%%%%%%%%%
%%%%%%%%%%%%%%%%%%%%%%%%%%%%%%%%%%%%%%
%%%%%%%%%%%%%%%%%%%%%%%%%%%%%%%%%%%%%%
%%%%%%%%%%%%%%%%%%%%%%%%%%%%%%%%%%%%%%
%%%%%%%%%%%%%%%%%%%%%%%%%%%%%%%%%%%%%%
%%%%%%%%%%%%%%%%%%%%%%%%%%%%%%%%%%%%%%
%%%%%%%%%%%%%%%%%%%%%%%%%%%%%%%%%%%%%%
%%%%%%%%%%%%%%%%%%%%%%%%%%%%%%%%%%%%%%



%%%%%%%%%%%%%%%%%%%%%%%%%%%%%%%%%%%%%%
%%%%%%%%%%%%%%%%%%%%%%%%%%%%%%%%%%%%%%
%%%%%%%%%%%%%%%%%%%%%%%%%%%%%%%%%%%%%%
%%%%%%%%%%%%%%%%%%%%%%%%%%%%%%%%%%%%%%
%%%%%%%%%%%%%%%%%%%%%%%%%%%%%%%%%%%%%%
%%%%%%%%%%%%%%%%%%%%%%%%%%%%%%%%%%%%%%
